% !TEX root = task-2-6-1.tex
% Editable task bank. Append new BankTask blocks at the end of the document.
\documentclass[12pt,a4paper]{article}
\usepackage{physics-textbook}
\usepackage{graphicx,needspace}
\usetikzlibrary{decorations.pathmorphing}
\newif\ifShowSolutions
\ifdefined\TutorEdition
  \ShowSolutionstrue
\else
  \ShowSolutionsfalse % Change to \ShowSolutionstrue to display solutions.
\fi
\newcounter{banktask}
\newsavebox{\BankTaskBox}
\newcommand{\BankTaskContent}[3]{%
  \noindent\textbf{\thebanktask.}\quad #2\par
  \ifShowSolutions
    \Needspace{50mm}% Keep the solution heading with its opening text or diagram.
    #3\par
    \ifdefined\TutorEdition
      \ifcsname TutorSource#1\endcsname
        \nopagebreak\smallskip{\small\csname TutorSource#1\endcsname\par}%
      \fi
    \fi
  \fi
}
\newcommand{\BankTask}[3]{%
  \refstepcounter{banktask}\par\medskip
  \sbox{\BankTaskBox}{\begin{minipage}{\linewidth}%
    \BankTaskContent{#1}{#2}{#3}\end{minipage}}%
  \ifdim\dimexpr\ht\BankTaskBox+\dp\BankTaskBox\relax>\dimexpr\textheight-20mm\relax
    \Needspace{35mm}\BankTaskContent{#1}{#2}{#3}%
  \else
    \Needspace{\dimexpr\ht\BankTaskBox+\dp\BankTaskBox\relax}%
    \noindent\usebox{\BankTaskBox}\par
  \fi
  \medskip
}
\begin{document}
\BookSubsection{2.6.1}{Термодинамический процесс}


% Source sdamgia-6065: https://phys-ege.sdamgia.ru/problem?id=6065
\BankTask{sdamgia-6065}{
Идеальный газ в количестве $\nu$ молей, имеющий концентрацию$n$ и находящийся при давлении $p$, сначала изобарически сжимают в $2$ раза, а затем изотермически расширяют в $4$ раза. Чему будут равны объем и температура этого газа в конце процесса расширения?

($k$ — постоянная Больцмана, $N_{A}$ — число Авогадро)

\TaskMatchingColumns{Величины}{А)  объем газа в конце процесса расширения\par Б)  температура газа в конце процесса\par расширения}{Их значения}{$1 ) \frac{\nu N_{\text{А}}}{2n}$\par $2 ) \frac{p}{2nk}$\par $3 ) \frac{2 \nu N_{\text{А}}}{n}$\par $4 ) \frac{2p}{nk}$}

\begin{center}\begin{tabular}{|c|c|}\hline
А & Б\\\hline
\strut & \strut\\\hline\end{tabular}\end{center}
}{
\smallskip\noindent\textbf{Решение}\par
Рассмотрим последовательно оба процесса.

$1$-$2$: изобарическое сжатие $p=const.$ В изобарическом процессе отношение $\frac{V}{T}$ остается постоянным. Следовательно, $T_2= \frac{V_2}{V_1} T_1= \frac{1}{2} T_1.$

$2$-$3$: изотермическое расширение $T=const.$ То есть $T_3=T_2= \frac{1}{2} T_1.$ В изотермическом процессе произведение $pV$ остается постоянным. Следовательно, $p_3=p_2 \frac{V_2}{V_3} = \frac{1}{4} p_1.$

Рассмотрим изменение объема в результате двух приведенных процессов: 

$V_2= \frac{1}{2} V_1,V_3=4V_2=2V_1.$

Концентрация $n= \frac{N}{V} ,$ следовательно, $n_3= \frac{1}{2} n_1.$ Теперь выразим искомые величины через величины, данные в условии задачи. Температура и давление идеального газа связаны соотношением $p=nkT.$ Откуда:

$T_3= \frac{p_3}{n_3k} = \frac{\frac{1}{4} p_1}{\frac{1}{2} n_1k} = \frac{p}{2nk}$ (Б — $2$).

Объем газа в конце процесса расширения найдем из уравнения Менделеева — Клапейрона:

$p_1V_1=\nu RT_1 \Longleftrightarrow V_1= \frac{\nu RT_1}{p_1} ,$

$V_3=2V_1= \frac{2\nu RT_1}{p_1} = \frac{2\nu R \frac{p_1}{n_1k}}{p_1} = \frac{2\nu N_{\text{А}}}{n}$ (А — $3$).

Ответ: $32$.
}

% Source sdamgia-2806: https://phys-ege.sdamgia.ru/problem?id=2806
\BankTask{sdamgia-2806}{
($N$ — число частиц, $p$ — давление, $V$ — объем, $T$ — абсолютная температура, $Q$ — количество теплоты)

\TaskMatchingColumns{Процессы}{А)  Изобарный процесс при $N = const$\par Б)  Изотермический процесс при $N = const$}{Формулы}{$1$)  $\frac{p}{T} = const$\par $2$)  $\frac{V}{T} = const$\par $3$)  $pV = const$\par $4$)  $Q = 0$}

\begin{center}\begin{tabular}{|c|c|}\hline
$A$ & Б\\\hline
\strut & \strut\\\hline\end{tabular}\end{center}
}{
\smallskip\noindent\textbf{Решение}\par
В изобарном процессе при неизменном количестве вещества $N=const$ — $\frac{V}{T} =const$ (А — $2$). В изотермическом процессе $pV=const$ (Б — $3$).

Ответ: $23$.
}


% Source E27EA7: https://ege.fipi.ru/bank/questions.php?search=1&proj=BA1F39653304A5B041B656915DC36B38&qid=E27EA7
\BankTask{E27EA7}{
При изучении процессов, происходящих с гелием, ученик занёс в таблицу результаты измерения температуры и давления одного и того же количества газа в различных равновесных состояниях. Газ считать идеальным.
\begin{center}
\begin{tabular}{|c|c|c|c|c|c|c|c|}
\hline
№ состояния & 1 & 2 & 3 & 4 & 5 & 6 & 7 \\
\hline
$p,\,\text{кПа}$ & $100$ & $90$ & $75$ & $50$ & $55$ & $75$ & $100$ \\
\hline
$t$, ${}^{\circ}\mathrm{C}$ & $27$ & $27$ & $27$ & $27$ & $57$ & $177$ & $327$ \\
\hline
\end{tabular}
\end{center}
\begin{enumerate}
\item Объём газа в состоянии $4$ в $2$ раза меньше объёма газа в состоянии $1$.
\item В состояниях $4$--$7$ объём газа был одинаковым.
\item Внутренняя энергия газа в состоянии $6$ в $3$ раза больше, чем в состоянии $5$.
\item При переходе от состояния $2$ к состоянию $3$ в ходе изотермического процесса газ получал тепло.
\item При переходе от состояния $5$ к состоянию $6$ в ходе изохорного процесса газ совершал работу.
\end{enumerate}
}{
\smallskip\noindent\textbf{Решение от GPT}\par
Для данного количества газа $V=\nu RT/p$, а для гелия $U=\frac32\nu RT$. Температуры состояний: $300$, $300$, $300$, $300$, $330$, $450$, $600\ \text{К}$.
\begin{enumerate}
\item Неверно: $V_4/V_1=p_1/p_4=2$, то есть объём больше в $2$ раза.
\item Верно: $T_4/p_4=T_5/p_5=T_6/p_6=T_7/p_7=6\ \text{К}/\text{кПа}$, поэтому объёмы одинаковы.
\item Неверно: $U_6/U_5=T_6/T_5=450/330=15/11\ne3$.
\item Верно: при постоянной температуре давление уменьшается, газ расширяется; $\Delta U=0$, $Q=A>0$.
\item Неверно: при изохорном процессе $A=0$.
\end{enumerate}
Ответ: $24$.
}

% Source 1F4F09: https://ege.fipi.ru/bank/questions.php?search=1&proj=BA1F39653304A5B041B656915DC36B38&qid=1F4F09
\BankTask{1F4F09}{
На рисунке показан график циклического процесса, проведённого с одноатомным идеальным газом, в координатах \(V\)–\(T\), где \(V\) — объём газа, \(T\) — абсолютная температура газа. Количество вещества газа постоянно.
\begin{center}
\begin{tikzpicture}[>=stealth,scale=0.95]
\draw[->] (0,0)--(4.7,0) node[below] {$T$};
\draw[->] (0,0)--(0,4.2) node[left] {$V$};
\node[below left] at (0,0) {$0$};
\coordinate (A) at (0.9,1.55);
\coordinate (B) at (1.8,3.1);
\coordinate (C) at (3.25,3.1);
\coordinate (D) at (0.9,0.86);
\draw[dashed,thin] (0,0)--(A) (0,0)--(D);
\draw (A)--(B)--(C)--(D)--(A);
\draw[->] (1.25,2.15)--(1.48,2.55);
\draw[->] (2.25,3.1)--(2.6,3.1);
\draw[->] (2.38,2.27)--(2.02,1.93);
\draw[->] (0.9,1.02)--(0.9,1.35);
\node[left] at (A) {$A$};
\node[above] at (B) {$B$};
\node[above right] at (C) {$C$};
\node[below right] at (D) {$D$};
\end{tikzpicture}
\end{center}
\begin{enumerate}
\item Давление газа в процессе \(CD\) постоянно, при этом над газом совершается положительная работа.
\item В процессе \(DA\) давление газа увеличивается.
\item В процессе \(AB\) газ получает положительное количество теплоты.
\item В состоянии \(D\) концентрация атомов газа максимальна.
\item В процессе \(BC\) внутренняя энергия газа уменьшается.
\end{enumerate}
}{
\smallskip\noindent\textbf{Решение от GPT}\par
Используем \(p=\nu RT/V\), \(U=\frac32\nu RT\), \(n=N/V\).
\begin{enumerate}
\item Верно: участок \(CD\) лежит на прямой, проходящей через начало координат, поэтому \(p=\mathrm{const}\). Объём уменьшается, работа над газом положительна.
\item Неверно: при \(T=\mathrm{const}\) объём на участке \(DA\) растёт, поэтому давление уменьшается.
\item Верно: на участке \(AB\) температура и объём увеличиваются, следовательно, \(\Delta U>0\), \(A>0\) и \(Q=\Delta U+A>0\).
\item Верно: в состоянии \(D\) объём минимален, поэтому концентрация максимальна.
\item Неверно: на участке \(BC\) температура растёт, поэтому внутренняя энергия увеличивается.
\end{enumerate}
\(\boxed{134}\).
}

% Source F5A444: https://ege.fipi.ru/bank/questions.php?search=1&proj=BA1F39653304A5B041B656915DC36B38&qid=F5A444
\BankTask{F5A444}{
Один моль одноатомного идеального газа участвует в процессе $1$--$2$--$3$, график которого изображён на рисунке в координатах $p$--$T$, где $p$~--- давление газа, $T$~--- абсолютная температура газа. Как изменяются объём газа $V$ в ходе процесса $1$--$2$ и плотность газа $\rho$ в ходе процесса $2$--$3$?

\begin{center}
\begin{tikzpicture}[x=1cm,y=1cm,>=stealth]
\draw[->] (0,0)--(4,0) node[below] {$T$};
\draw[->] (0,0)--(0,3.3) node[left] {$p$};
\node[below left] at (0,0) {$0$};
\draw (0,0)--(2.9,1.22)--(2.9,3.05);
\fill (1.03,0.433) circle (1.2pt) node[above left] {$1$};
\fill (2.9,1.22) circle (1.2pt) node[right] {$2$};
\fill (2.9,3.05) circle (1.2pt) node[right] {$3$};
\draw[->] (1.4,0.589)--(1.95,0.82);
\draw[->] (2.9,1.7)--(2.9,2.2);
\end{tikzpicture}
\end{center}
Для каждой величины определите соответствующий характер изменения:

\begin{enumerate}
\item увеличивается
\item уменьшается
\item не изменяется
\end{enumerate}


\TaskChangeTable{Объём газа в ходе процесса $1$--$2$}{Плотность газа в ходе процесса $2$--$3$}
}{
\smallskip\noindent\textbf{Решение от GPT}\par
Из $pV=\nu RT$ следует $V=\nu RT/p$. На участке $1$--$2$ прямая проходит через начало координат, поэтому $p/T=\mathrm{const}$ и объём не изменяется (3).

На участке $2$--$3$ температура постоянна, давление увеличивается, поэтому объём уменьшается. При постоянной массе $\rho=m/V$ увеличивается (1).

Ответ: $31$.
}

% Source F0D87A: https://ege.fipi.ru/bank/questions.php?search=1&proj=BA1F39653304A5B041B656915DC36B38&qid=F0D87A
\BankTask{F0D87A}{
$1$ моль одноатомного идеального газа участвует в процессе $1$--$2$--$3$, график которого изображён на рисунке в координатах $p$--$V$, где $p$~--- давление газа, $V$~--- объём газа. Как изменяются концентрация молекул $n$ газа в ходе процесса $1$--$2$ и абсолютная температура $T$ газа в ходе процесса $2$--$3$? Масса газа остаётся постоянной.

\begin{center}
\begin{tikzpicture}[x=1cm,y=1cm,>=stealth]
\draw[->] (0,0)--(4.2,0) node[below] {$V$};
\draw[->] (0,0)--(0,3.5) node[left] {$p$};
\node[below left] at (0,0) {$0$};
\draw (0.8,0.9)--(3.2,0.9)--(3.2,2.9);
\fill (0.8,0.9) circle (1.2pt) node[below left] {$1$};
\fill (3.2,0.9) circle (1.2pt) node[below right] {$2$};
\fill (3.2,2.9) circle (1.2pt) node[above right] {$3$};
\draw[->] (1.4,0.9)--(2,0.9);
\draw[->] (3.2,1.4)--(3.2,1.95);
\end{tikzpicture}
\end{center}
Для каждой величины определите соответствующий характер изменения:

\begin{enumerate}
\item увеличивается
\item уменьшается
\item не изменяется
\end{enumerate}

\TaskChangeTable{Концентрация молекул газа в ходе процесса $1$--$2$}{Абсолютная температура газа в ходе процесса $2$--$3$}
}{
\smallskip\noindent\textbf{Решение от GPT}\par
На участке $1$--$2$ объём увеличивается. Число молекул $N$ постоянно, поэтому концентрация $n=N/V$ уменьшается (2).

На участке $2$--$3$ объём постоянен, давление увеличивается. Из $T=pV/(\nu R)$ следует, что температура увеличивается (1).

Ответ: $21$.
}

% Source F12DcD: https://ege.fipi.ru/bank/questions.php?search=1&proj=BA1F39653304A5B041B656915DC36B38&qid=F12DcD
\BankTask{F12DcD}{
\begin{center}
\begin{tikzpicture}[>=stealth,scale=0.85]
\draw[->] (0,0)--(4,0) node[below] {$T$};
\draw[->] (0,0)--(0,3.8) node[left] {$V$};
\node[below left] at (0,0) {$0$};
\coordinate (A) at (2.5,3.2);
\coordinate (B) at (2.5,1);
\coordinate (C) at (1.3,0.52);
\draw[dashed] (0,0)--(C);
\draw (A)--(B)--(C);
\draw[->] (2.5,2.65)--(2.5,1.95);
\draw[->] (2.25,0.9)--(1.75,0.7);
\fill (A) circle (1.5pt) node[right] {$1$};
\fill (B) circle (1.5pt) node[right] {$2$};
\fill (C) circle (1.5pt) node[above left] {$3$};
\end{tikzpicture}
\end{center}
Один моль идеального газа участвует в процессе \(1\text{–}2\text{–}3\), график которого изображён на рисунке в координатах \(V\text{–}T\), где \(V\) — объём газа, \(T\) — абсолютная температура газа. Как изменяются давление газа \(p\) в ходе процесса \(1\text{–}2\) и концентрация молекул газа \(n\) в ходе процесса \(2\text{–}3\)? Масса газа остаётся постоянной.

Для каждой величины определите соответствующий характер изменения:
\begin{enumerate}
\item увеличивается
\item уменьшается
\item не изменяется
\end{enumerate}
\TaskChangeTable{Давление газа в ходе процесса \(1\text{–}2\)}{Концентрация молекул газа в ходе процесса \(2\text{–}3\)}
}{
\smallskip\noindent\textbf{Решение от GPT}\par
В процессе \(1\text{–}2\) температура постоянна, а объём уменьшается. Из \(p=\nu RT/V\) следует, что давление увеличивается (\(1\)). В процессе \(2\text{–}3\) объём уменьшается при постоянном числе молекул, поэтому концентрация \(n=N/V\) увеличивается (\(1\)).

Ответ: \(11\).
}

\end{document}
